\section{Conclusion}

This paper has created a prescriptive decision-support model to optimize the sales of third-party delivery platforms, based on the data provided by Wolt and Foodora. The findings revealed that decision making can be greatly enhanced through integration of data, the application of forecasting models and the application of optimization techniques. 

The forecasting model was able to determine the demand patterns of weekly seasonality and changes in the order volume. These forecasts were converted into actionable recommendations by the optimization model, and it showed the possibility of optimizing the profit by increasing the platform-level allocation. 

In general, the research establishes the fact that predictive and prescriptive analytics could be combined to offer a solid basis to data-oriented decision-making in the QSR industry. 

This research adds to the scholarly community in that it combines the forecasting and optimization into a Design Science Research model. It builds upon the existing research by integrating predictive analytics with prescriptive decision-making on the platform of third-party delivery. The gap in the multi-platform analysis in the data of food delivery is also discussed in the study. 

Practically, the research offers a systematic framework that can be applied by companies like HungryHub Oy to enhance decision-making process. The framework assists in converting raw sales information into actionable insights to aid in enhancing planning, allocation of resources, and optimization of profits across delivery platforms.


\section{Limitations and Future Work}

\subsection{Limitations}

The study has a number of limitations in spite of its contributions. First, the analysis is conducted with the help of a small set of data of two platforms only, Wolt and Foodora. Second, the model is based on platform-specific assumptions, including fixed commission structure and simplified constraints. Third, the framework is not subjected to real-time operational environment and is tested with historical data.


\subsection{Future Work}

There are several possible ways in which this research can be extended in the future. One direction you can take is in adding multiple restaurants or business units to improve generalization. The other alternative is to develop real time decision-supporting systems that can update their future forecasts and recommendations in real time.

Additionally, future research could include advanced optimization techniques, dealing with operational variables like staffing, inventory, and logistics for the implementation of the delivery.